\section{Experiments}
\label{sec:experiments}

\subsection{Questions and Frozen Comparisons}
We ask whether iterative diffusion improves over a one-step estimator (H1),
whether matching support improves over population and same-cell wrong support
(H2), and whether natural-EEG preservation remains acceptable (H3).  The core
methods are RAW, an existing population deterministic cleaner (POP), DIFF-POP,
DIFF-MATCH, DIFF-WRONG, and DET-MATCH.  The three diffusion arms use the same
checkpoint and common random numbers.  We additionally report $K=1$ versus
$K=8$ and removal of identity examples as targeted ablations.

All models use three fixed training seeds (20260811--20260813), at most 12,000
updates, independent validation-optimal checkpoints, and the same window and
normalization pipeline.  Diffusion uses a cosine schedule, $v$-prediction, EMA,
Min-SNR weighting, DDIM25, and an arithmetic mean of eight samples.  Statistical
units are source records or participant stems, never windows.  Windows and seeds
are first averaged within a unit; 20,000 paired bootstrap samples then provide
descriptive 95\% intervals.  SGE bootstrap sampling is stratified by exact
montage--reference--sampling-rate cell.

\subsection{Datasets and Evidence Scope}
\paragraph{Klados--Bamidis v4.}
We use the frozen split with source records 1--30 for training, 31--36 and
44--45 for development, and all 16 records 37--43 and 46--54 for evaluation.
The recordings combine real EEG and real ocular sources with a common clean
target.  Participant provenance is not reliable, so this is paired
source-record mechanism evidence, not participant-held-out evidence.  The
primary endpoint is clean RRMSE; correlation, artifact reconstruction error,
$\Delta$SNR, and complement error are secondary.

\paragraph{SGEYESUB.}
All five studies are used as real-EEG development data.  For each participant
stem, block 1 is support and block 2 is query.  Operators and wrong controls
never cross exact compatibility cells.  Query EOG and annotations are opened
only after all method outputs are frozen.  Of 59 available stems, 58 are
compatible; study05/study05\_p42 lacks a compatible population cell and remains
in the coverage denominator.  The primary artifact endpoint is independent EOG
coherence reduction.  Safety endpoints are low-artifact observation
preservation, ERP preservation, continuous-segment Welch PSD distortion,
reference-free covariance distortion, output scale, and observation change.
No clean-waveform claim is made for SGEYESUB.

\paragraph{Independent participant data.}
The official EEGEyeNet OSF page redirects its synchronized release to a public
Google Drive.  Two scheduled, resumable attempts reached the participant MAT
folders but the official server denied public access to participant files.
Consequently, no independent EEGEyeNet outcome was opened and the current
experiment remains development evidence.

\subsection{Technical Validity}
The implementation passed bounded-target overfit, zero-correction identity,
finite reverse trajectory, coefficient bounds, complement consistency
($\le10^{-5}$), operator intervention, EMA checkpoint reload/resume, explicit
CUDA generator state, participant-specific seeds, and forbidden-query-field
tests.  Training completed all 78 dataset--fold--seed tasks.  A first Klados
evaluation attempt exposed an FP64/FP32 mismatch in the diagnostic complement
metric; only that post-processing conversion was corrected, and all three
Klados inference seeds were replayed.

\subsection{Paired Waveform Results (H1)}
Table~\ref{tab:klados_results} reports source-record means.  Lower RRMSE is
better.  DIFF-MATCH improves over the information-matched DET-MATCH by a utility
effect of $0.0592$ (95\% CI $[0.0336,0.0879]$, $n=16$), supporting H1 under the
tested protocol.  POP remains stronger than either learned matching model; H1
therefore identifies an iterative-versus-one-step difference, not superiority
over every baseline.

\begin{table}[t]
\caption{Klados paired mechanism results (three-seed averages).}
\label{tab:klados_results}
\centering
\begin{tabular}{lrrr}
\toprule
Method & Clean RRMSE $\downarrow$ & Correlation $\uparrow$ & $\Delta$SNR dB $\uparrow$\\
\midrule
RAW & 1.2948 & 0.6302 & 0.000\\
POP & 0.3990 & 0.9204 & 10.402\\
DIFF-POP & 0.4767 & 0.8932 & 8.925\\
DIFF-MATCH & 0.5325 & 0.8744 & 7.940\\
DIFF-WRONG & 0.5543 & 0.8641 & 7.751\\
DET-MATCH & 0.5917 & 0.8497 & 7.017\\
\bottomrule
\end{tabular}
\end{table}

The eight-sample posterior also supplies useful uncertainty.  Its
uncertainty--error correlation is 0.445 versus 0.140 for a three-seed
deterministic ensemble, and its risk--coverage AUC is lower (0.241 versus
0.292).  These are development diagnostics rather than an independently
confirmed clinical risk model.

\subsection{Natural EEG Subject Intervention and Safety (H2--H3)}
Table~\ref{tab:sge_results} shows natural-EEG means.  Matching support does not
improve the independent EOG endpoint over population:
$\Delta_{\mathrm{match-pop}}=-0.00515$
($[-0.0113,0.00071]$, $n=58$).  Its advantage over a same-cell wrong operator is
$0.00286$ ($[-0.00303,0.00867]$), also inconclusive.  H2 is therefore not
supported.

\begin{table}[t]
\caption{SGEYESUB real-EEG development results.  EOG reduction and preservation
are higher-is-better; PSD distortion is lower-is-better.}
\label{tab:sge_results}
\centering
\begin{tabular}{lrrr}
\toprule
Method & EOG reduction & Preservation & PSD distortion\\
\midrule
RAW & 0.0000 & 1.0000 & 0.0000\\
POP & 0.3231 & 0.7762 & 0.3830\\
DIFF-POP & 0.2375 & 0.8973 & 0.1509\\
DIFF-MATCH & 0.2323 & 0.8891 & 0.1602\\
DIFF-WRONG & 0.2295 & 0.8938 & 0.1543\\
DET-MATCH & 0.2301 & 0.8683 & 0.2039\\
\bottomrule
\end{tabular}
\end{table}

Against the frozen $-0.02$ population-relative margin, DIFF-MATCH passes all
four safety endpoints: low-artifact preservation, ERP preservation,
continuous-segment Welch PSD distortion, and covariance distortion.  Identity
examples materially improve safety: removing them reduces preservation from
0.889 to 0.796 and increases PSD distortion from 0.160 to 0.358.  H3 is
supported in development.

\subsection{Interpretation}
The complete result is mixed and constrains the paper's claim.  Iterative
coefficient diffusion improves over an information-matched one-step model and
its posterior variance provides better error ranking.  The structural operator
conditioning is numerically active and safe, but the matching support operator
does not outperform population or same-cell wrong support at the participant
level.  Thus the current experiment supports diffusion in an estimated artifact
subspace, while the stronger subject-aware transfer claim remains unverified.
